% ===================================================================
% CHAPTER 10: FUTURE WORK
% ===================================================================
\chapter{Идни планови и развојна патека}
\label{chap:future}

Насоките подолу се подредени по приоритет: прво она што директно произлегува од измерените
наоди, потоа она што ги проширува границите на генерализација.

\section{Калибрација: од оракул праг кон праг употреблив во распоредување}

Преминот низ праговите веќе покажа дека колапсот е калибрациски и дека со подесен праг SISA
достигнува избалансирана точност 0,709 наспроти 0,671 за наивниот пристап
(\hyperref[chap:results]{Дел~\ref*{chap:results}}, §7.5). Она што останува не е да се утврди
\emph{дали} калибрацијата помага, туку \textbf{колку од тоа преживува во реални услови}.

\begin{itemize}
  \item \textbf{Праг избран надвор од множеството на оценување.} Тековната вредност е оракул: прагот
  ја максимизира избалансираната точност врз истото множество врз кое се известува
  (\hyperref[chap:limitations]{Дел~\ref*{chap:limitations}}, §8.13). Соодветната постапка го
  бира прагот на издвоено валидациско множество и го оценува на дисјунктно тест-множество.
  Разликата меѓу двете вредности е вистинската цена на калибрацијата и е директно изводлива
  од постојните артефакти.
  \item \textbf{Стабилност на прагот меѓу извршувања.} Оптималниот праг варира (0,947--0,987 за
  SISA, 0,964--0,996 за наивниот пристап). Ако варира и меѓу распоредувања, фиксен праг нема
  да биде доволен и ќе биде потребна периодична рекалибрација --- што е оперативен трошок што
  овој труд не го мери.
  \item \textbf{Калибрација на веројатностите} (Platt-скалирање, изотона регресија) наместо само
  избор на праг, што дава и толкувливи веројатности. Особено релевантно бидејќи \textbf{двата}
  крака бараат екстремни прагови (0,96--0,99), што укажува на системска лоша калибрација, не
  само на поместена граница кај SISA.
  \item \textbf{Оценување со тест-множество по почетна вредност}, за да се отстрани преостанатата
  оптимистичка пристрасност од заедничкото тест-множество
  (\hyperref[chap:limitations]{Дел~\ref*{chap:limitations}}, §8.13).
\end{itemize}

\section{Алтернативна агрегација на составните модели}

Ако калибрацијата не го затвори јазот, причината треба да се нападне директно. Колапсот е
припишан на \textbf{параметарското просечување} на пет независно обучени мрежи --- отстапување од
класичниот SISA направено за клиентот воопшто да може да учествува во FedAvg.

Можни насоки:

\begin{itemize}
  \item \textbf{Ансамбл од предвидувања на ниво на клиент.} Клиентот ги задржува петте составни
  модели и ги спојува нивните излези локално, а кон серверот испраќа единствен модел
  дестилиран од тој ансамбл. Ова го задржува оригиналниот механизам на SISA, по цена на чекор
  на дестилација.
  \item \textbf{Порамнување на невроните пред просечувањето.} Просечувањето тежини на независно
  обучени мрежи страда од тоа што нивните скриени единици не се во соодветен распоред.
  Порамнувањето на пермутациите пред просечувањето е познат пристап што вреди да се испита во
  оваа поставка.
  \item \textbf{Помалку шардови, повеќе сегменти.} Изборот $S = 5$ ја определува и предноста при
  обновувањето и бројот мрежи што се просечуваат. Испитување на просторот $S \times R$ би
  покажало дали постои точка во која предноста во цената се задржува, а колапсот исчезнува.
\end{itemize}

\section{Довршување на претходно регистрираните помошни извршувања}

Три серии од претходната регистрација останаа неизвршени и се директно наведени како
ограничување во \hyperref[chap:limitations]{Дел~\ref*{chap:limitations}}:

\begin{itemize}
  \item \textbf{Чисти референтни извршувања} (\texttt{POISON\_MODE=off}), по едно на почетна
  вредност. Тие ја даваат горната граница на корисноста и се единствениот начин хипотезата H5
  да се оцени во нејзината првично регистрирана форма.
  \item \textbf{Студија на чувствителност со скалиран модел.} Наодот за H4 --- дека
  влезно-излезните операции се занемарливи --- важи за модел со 12\,865 параметри.
  Повторување на споредбата со модел за еден или два реда на големина поголем би ја лоцирала
  точката во која записите на контролните точки стануваат доминантниот трошок на SISA. Ова е
  најинформативната единечна дополнителна серија, бидејќи токму тоа е претпоставката врз која
  се заснова изборот на SD-картичка од класата A1 како тесно грло.
  \item \textbf{Извршување под WAN-услови.} Мерниот прозорец е чисто локална пресметка, па не се
  очекува влијание --- но тоа е очекување, не измерен факт.
\end{itemize}

Дополнително, \textbf{преправање на четирите пара извршени во различни сесии} во иста сесија би ја
отстранило единствената преостаната методолошка забелешка, иако проверката на робусноста веќе
покажува дека таа не го менува резултатот.

\section{Проширување на опсегот на генерализација}

\begin{itemize}
  \item \textbf{Повеќе уреди и повеќе картички.} Тековниот труд е студија на случај врз еден
  примерок. Повторување врз неколку уреди и неколку SD-картички би ја одвоило варијабилноста
  на примерокот од ефектот на методот, и би дозволило тврдење за класата уреди, наместо за
  еден уред.
  \item \textbf{Контролирана амбиентална околина.} Најголемото методолошко ограничување е
  незапишаната амбиентална температура. Мерења во просторија со контролирана температура ---
  или барем со запишан сензор --- би ги направиле апсолутните термички вредности споредливи
  меѓу извршувања.
  \item \textbf{Поголема федерација.} Со четири клиенти, придонесот на компромитираниот јазол во
  глобалниот модел е околу една четвртина. Со десетици јазли, разредувањето е поголемо и
  прашањето за приближното глобално заборавање добива друга тежина.
  \item \textbf{Други податочни множества и избалансирани задачи.} Колапсот во мнозинската класа
  најверојатно зависи од нерамнотежата од 96,5\,\%. Повторување врз поизбалансирана задача би
  покажало дали параметарското просечување е проблематично воопшто, или само при изразена
  нерамнотежа.
\end{itemize}

\section{Од оракул кон вистинско откривање}

Во целиот труд, компромитираните примери се дадени преку оракул-маска; откривањето е
експлицитно надвор од опсегот. Во реален систем тоа е сам по себе тежок проблем, а неговата
цена --- пресметковна и временска --- не е вклучена во ниту еден од измерените трошоци.

Природното продолжение е да се спои постапката за заборавање со механизам за откривање и да
се измери \textbf{целосната} патека од сомневање до обновен модел. Тоа го отвора и прашањето за
цената на лажните аларми: заборавање активирано по погрешна детекција троши исти ресурси, а
уништува валидни податоци.

\section{Заборавање на ниво на клиент}

Овој труд мери заборавање на \textbf{ниво на примероци} кај еден јазол. Регулаторно почестиот
случај е повлекување на \textbf{целиот} јазол од федерацијата. Таа поставка ја менува
пресметката суштински --- SISA на ниво на клиент повеќе не е локална постапка, туку прашање за
реконструкција на глобалниот модел, каде пристапите што чуваат историја на ажурирања
\cite{liu2021federaser, cao2023fedrecover} стануваат директни конкуренти. Споредба на тие
пристапи по истата методологија на физичко мерење би била природното продолжение на овој
труд.
